\scanpage{55}
Alors $\varinjlim_i|X_i|$ est un espace loc. compact dénombrable à l'$\infty$, et les $|X_i|$ forment un système fond. de compacts. Le foncteur
\[
\mathfrak X\longmapsto|\mathfrak X|=\varinjlim_i|X_i|
\]
de la catégorie \inserted{$\widetilde{\mathcal M}$} des ind-objets spéciaux \inserted{de $\mathcal M$} vers les espaces loc. compacts est exact à g (\uline{NB} $\mathcal M$ est stable par $\varprojlim$ finies dans $\widehat{\mathcal M}$), conservatif, peut-être pas transportable -- ce qui nous permet d'introduire la notion de structure \uline{pseudo} modérée sur un espace loc. compact (s'il est compact, cela revient à la structure modérée), de sorte que la catégorie des ind-objets spéciaux devient équivalente à celle des espaces loc. cpt pseudo modérés. On définit encore la notion de \uline{morphismes \inserted{pseudo} modérés}, de sous-ensemble (fermé) pseudo-modéré d'un tel $X$. On a que si $f:X-Y$ est \uline{propre} et pseudo-modéré, $f(X)$ est une partie pseudo-modérée de $Y$.
\margin{\unsure{compt.}}

\uline{Donnée supplémentaire} a) Structure pseudo-modérée sur $\mathbb R$ pour laquelle les intervalles $I_n=[-n,+n]$ $(n\in\mathbb Z)$ sont des sous-ens. \inserted{(ps)} modérés, la multiplication par \deleted{$\unsure{n\in\mathbb N^*}$}\inserted{$n\in\mathbb Z$} étant un morphisme (donc isomorphisme \inserted{si $n\neq0$}) ps modérés \unsure{$g$} de $\mathbb R$. On aura donc
\[
\mathbb R=``\varinjlim'' I_n\simeq\varinjlim I_1,
\]
les morphismes de transition $I_1(n)-I_1(n+1)$ étant la multiplication par $\dfrac n{n+1}$.
Donc la donnée a) revient à la donnée b) Structure modérée sur $I_1=[-1,+1]$ pour laquelle les applications
\[
t\longmapsto rt\qquad(r\in\mathbb Q^*,\;\unsure{\text{ayant}}\;|r|\leq1)
\]
sont modérées. On suppose de plus

M4) $\left\{\begin{array}{l}
1^\circ)\text{ le point }1\text{ de }I_1\text{ (ou de }\mathbb R,\text{ c'est kif kif) est modéré}\\
2^\circ)\text{ l'application somme }\mathbb R\times\mathbb R\xrightarrow{x+y}\mathbb R\\
\qquad(\text{i.e. }(x,y)\longmapsto\dfrac{x+y}{2}:\;I_1\times I_1-I_1,\\
\qquad\text{c'est kif-kif) est modérée}\\
3^\circ)\text{ l'application produit }\mathbb R\times\mathbb R\xrightarrow{xy}\mathbb R\\
\qquad(\text{i.e. }(x,y)\longmapsto xy:\;I_1\times I_1-I_1,\\
\qquad\text{c'est kif-kif) est modérée}
\end{array}\right.$

de $1^\circ$ et $2^\circ$, on conclut que tous les points $x\in\mathbb Q\subset\mathbb R$ sont modérés, que les translations \deleted{\unsure{et les homothéties}} $\tau_{\unsure{x|n}}$ \inserted{$x\in\mathbb Q$} sont modérées, que \unsure{tous intervalles|tout intervalle} de $\mathbb R$ à extrémités modérées sont \unsure{itou}.

\margin{56}
